\chapter{Future Work}
\label{ch:5}
\noindent

A central limitation of the present framework is that it does not learn fragmentation itself.
The model operates on \emph{precomputed} derivation graphs.
Candidate fragments and their parent--child relations are generated by the rule-based graph-transformation pipeline upstream, and are then encoded by the neural model.
The learning problem thus begins only \emph{after} a fragmentation space has been constructed.
The network learns how to embed, combine, align, and score molecular, fragment, and spectral information, but not which fragmentation events to propose in the first place.
This design is appropriate for a first proof of concept, since it keeps the fragment-level representation explicit and the learning problem tractable, but it also marks the clearest boundary of the current approach and thus the most important direction for future work.

\section{From Fragment Encoding to Fragment Generation}
The present implementation is best understood as a \emph{fragment-aware encoder} rather than a full fragmentation simulator.
Its fragment branch consumes fragment graphs together with their derivation structure encoding a graph whose nodes are fragment graphs and whose directed edges represent parent--child relations extracted from the precomputed derivation graph.
The model therefore receives a structured fragmentation hypothesis as input instead of inferring it autonomously.
The natural next step is to move from \emph{encoding} fragmentation structures to \emph{generating} them: rather than assuming the relevant fragment space has already been enumerated, the model would be trained to predict which bonds or substructures are most likely to fragment, in which order, and under which contextual conditions.
This would shift the system from symbolic preprocessing with neural downstream learning toward a more integrated neuro-symbolic architecture, considerably closer to a true mechanistic simulator of fragmentation pathways.


\section{Toward a Generate--Score Architecture}
A key source of inspiration for this transition is ICEBERG \cite{Goldman2024}, which separates the task into a \emph{Generate} module that predicts likely fragmentation events and a \emph{Score} module that assigns intensities to the resulting fragments.
Crucially, the fragment graph is not treated as fixed input, but rather as a latent object that must be predicted itself.
The generative stage is trained to reconstruct an estimated fragmentation hypergraph, while a second module predicts intensities and hydrogen rearrangements of the generated fragments.
The present model already supplies the main ingredients for the \emph{Score} side of such a system.
It embeds molecular graphs, fragment-level structures, and spectra into a shared latent space; it predicts spectra from latent representations; and it uses bidirectional training signals such as contrastive alignment and reconstruction losses.
What is missing is an upstream module that decides which fragments enter the fragment branch at all.
The symbolic derivation graphs produced by the rule-based pipeline could be used to supervise such a generator — initially as a learnt pruning mechanism over the derivation graph and later as a more ambitious autoregressive model that predicts fragmentation events directly (\Cref{fig:generative-future}).


\begin{figure}[htbp]
    \centering
    \resizebox{0.9\textwidth}{!}{%
        \begin{tikzpicture}[
    fig,
    bigarrow/.style={-{Latex[length=3.5mm,width=3.2mm]}, fig heavy},
    greenA/.style={bigarrow, SeaGreen},
    blueA/.style={bigarrow, RoyalBlue},
    grayA/.style={bigarrow, black!55},
    treearrow/.style={-{Latex[length=1.6mm]}, dashed, black!60, thin},
    fragnode/.style={circle, draw=SeaGreen!80, fill=SeaGreen!50, minimum size=4mm, inner sep=0pt},
    fragempty/.style={circle, draw=SeaGreen!70, fill=white, minimum size=4mm, inner sep=0pt},
    fraglight/.style={circle, draw=SeaGreen!60, fill=SeaGreen!25, minimum size=4mm, inner sep=0pt},
    fragroot/.style={circle, draw=SeaGreen!80, fill=SeaGreen!70, minimum size=5mm, inner sep=0pt},
    grayN/.style={circle, draw=black!41, fill=black!22.55, minimum size=2.6mm, inner sep=0pt},
]


% =================== Color palette ===================


% =================================================================
%   STAGE A: Precursor molecule box
% =================================================================
\begin{scope}[shift={(6, 18.60)}]
    \draw[rounded corners=2pt, black, thick]
        (-5.0,-1.05) rectangle (5.0, 1.05);
    \node[anchor=south, font=\Large\bfseries] at (0, 1.20) {Precursor molecule};

    % --- Chemical structure ---
    % Benzene ring centered at (-2.7, 0), R=0.55
    \def\bcx{-3.2}
    \def\bcy{0}
    \def\bR{0.55}
    \foreach \i in {0,...,5} {
        \pgfmathsetmacro{\ang}{60*\i}
        \coordinate (b\i) at ($ (\bcx,\bcy) + (\ang:\bR) $);
    }
    % Hexagon edges (alternating double bonds for aromatic look)
    \draw[thick] (b0)--(b1) (b2)--(b3) (b4)--(b5);
    \draw[thick, double distance=1.2pt] (b1)--(b2) (b3)--(b4) (b5)--(b0);

    % Backbone zigzag from benzene right vertex to N+
    \def\zy{0.20}        % up vertex y
    \def\zd{0.0}         % down vertex y baseline
    \def\dx{0.55}        % step
    % Start at b0 (right vertex of benzene at angle 0)
    \coordinate (bb0) at (b0);
    \coordinate (bb1) at ($ (b0) + (\dx, \zy) $);          % CH2 (implicit C)
    \coordinate (bb2) at ($ (bb1) + (\dx,-\zy) $);         % NH
    \coordinate (bb3) at ($ (bb2) + (\dx, \zy) $);         % C(=O)
    \coordinate (bb4) at ($ (bb3) + (\dx,-\zy) $);         % NH
    \coordinate (bb5) at ($ (bb4) + (\dx, \zy) $);         % CH2
    \coordinate (bb6) at ($ (bb5) + (\dx,-\zy) $);         % CH2
    \coordinate (bb7) at ($ (bb6) + (\dx, \zy) $);         % N+
    \coordinate (oxy) at ($ (bb3) + (0, 0.55) $);          % carbonyl O

    % Bonds (skeletal) — leave gaps near labeled atoms
    \draw[thick] (bb0)--(bb1);
    \draw[thick] (bb1)--($ (bb2) + (-0.10, 0.04) $);
    \draw[thick] ($ (bb2) + ( 0.10, 0.04) $)--(bb3);
    \draw[thick, double distance=1.2pt] (bb3)--($ (oxy) + (0,-0.13) $);
    \draw[thick] (bb3)--($ (bb4) + (-0.10,-0.04) $);
    \draw[thick] ($ (bb4) + ( 0.10,-0.04) $)--(bb5);
    \draw[thick] (bb5)--(bb6);
    \draw[thick] (bb6)--($ (bb7) + (-0.10,-0.04) $);

    % Atom labels
    \node[font=\small, Salmon] at (oxy) {O};
    \node[font=\small] at (bb2) {N};
    \node[font=\scriptsize] at ($ (bb2) + (0,-0.32) $) {H};
    \draw[thick] ($ (bb2) + (0,-0.10) $) -- ($ (bb2) + (0,-0.22) $);
    \node[font=\small] at (bb4) {N};
    \node[font=\scriptsize] at ($ (bb4) + (0,-0.32) $) {H};
    \draw[thick] ($ (bb4) + (0,-0.10) $) -- ($ (bb4) + (0,-0.22) $);
    \node[font=\small] at (bb7) {N};
    \node[font=\tiny] at ($ (bb7) + (0.20, 0.18) $) {$+$};

    % Three methyl groups on N+
    \draw[thick] ($ (bb7) + (0.12, 0.04) $) -- ($ (bb7) + (0.55, 0.30) $);   % up-right
    \draw[thick] ($ (bb7) + (0.10,-0.10) $) -- ($ (bb7) + (0.55,-0.40) $);   % down-right
    \draw[thick] ($ (bb7) + (0.00,-0.16) $) -- ($ (bb7) + (0.00,-0.55) $);   % down
\end{scope}

% Arrow from precursor to Generate
\draw[greenA] (6, 17.45) -- (6, 16.55);

% =================================================================
%   STAGE B: GENERATE module (green)
% =================================================================
\begin{scope}[shift={(6, 12.10)}]
    % Outer box
    \draw[rounded corners=3pt, SeaGreen, fig medium, fill=SeaGreen!10]
        (-4.8,-2.00) rectangle (4.8, 4.45);
    % Title
    \node[anchor=north, font=\LARGE\bfseries, text=SeaGreen] at (0, 4.30) {Generate};
    \node[anchor=north, font=\small\itshape] at (0, 3.65) {(neural fragmentation module)};

    % Cloud icon (top-left of module) — proper cloud silhouette + small graph
    \begin{scope}[shift={(-3.4, 3.7)}]
        % Cloud silhouette (TikZ cloud shape)
        \node[cloud, draw=SeaGreen, fig medium, fill=white,
              cloud puffs=11, cloud puff arc=135, aspect=1.6,
              minimum width=1.55cm, minimum height=0.95cm,
              inner sep=0pt] (cloudbg) at (0,0) {};
        % Small graph inside (4 nodes, triangle + extra)
        \node[circle, fill=SeaGreen, minimum size=2mm, inner sep=0pt] (cd1) at (-0.20, -0.05) {};
        \node[circle, fill=SeaGreen, minimum size=2mm, inner sep=0pt] (cd2) at ( 0.20, -0.05) {};
        \node[circle, fill=SeaGreen, minimum size=2mm, inner sep=0pt] (cd3) at ( 0.00,  0.18) {};
        \node[circle, fill=SeaGreen, minimum size=2mm, inner sep=0pt] (cd4) at ( 0.00, -0.20) {};
        \draw[SeaGreen, fig thin] (cd1)--(cd3) (cd2)--(cd3)
                                          (cd1)--(cd4) (cd2)--(cd4)
                                          (cd1)--(cd2);
    \end{scope}

    % Inner sub-box
    \draw[rounded corners=2pt, SeaGreen!60, thin]
        (-4.5,-1.70) rectangle (4.5, 2.55);
    \node[anchor=north, font=\small, align=center] at (0, 2.40)
        {Predict fragmentation events and candidate fragments};

    % --- Fragmentation tree ---
    % Root at top-center
    \node[fragroot] (rt) at (0.5, 1.25) {};
    % Level 2 (3 children)
    \node[fragnode] (a1) at (-1.4, 0.20) {};
    \node[fragnode] (a2) at ( 0.5, 0.20) {};
    \node[fragnode] (a3) at ( 2.4, 0.20) {};
    % Edges
    \foreach \c in {a1, a2, a3} { \draw[treearrow] (rt) -- (\c); }
    % Level 3 (children)
    \node[fraglight] (l1) at (-2.4,-1.10) {};
    \node[fraglight] (l2) at (-1.5,-1.10) {};
    \node[fraglight] (l3) at (-0.6,-1.10) {};
    \node[fraglight] (l4) at ( 0.3,-1.10) {};
    \node[fraglight] (l5) at ( 1.2,-1.10) {};
    \node[fragempty] (l6) at ( 2.1,-1.10) {};
    \node[fragempty] (l7) at ( 3.0,-1.10) {};
    \node[font=\small] at ( 3.65,-1.10) {$\dots$};
    % Edges from level 2 to level 3
    \draw[treearrow] (a1) -- (l1);
    \draw[treearrow] (a1) -- (l2);
    \draw[treearrow] (a2) -- (l3);
    \draw[treearrow] (a2) -- (l4);
    \draw[treearrow] (a2) -- (l5);
    \draw[treearrow] (a3) -- (l6);
    \draw[treearrow] (a3) -- (l7);
\end{scope}

% =================================================================
%   STAGE C: MØD pipeline box (left side)
% =================================================================
\begin{scope}[shift={(-1.0, 12.10)}]
    \draw[rounded corners=3pt, black!60, fill=black!12]
        (-1.5,-2) rectangle (1.5, 2.0);
    % Mini tree icon at top — same shape as fragmentation DAG/tree but in gray
    \begin{scope}[shift={(0, 1.10)}]
        % Root
        \node[circle, draw=black!41, fill=black!28.7,
              minimum size=3.4mm, inner sep=0pt] (gr) at (0, 0.40) {};
        % Children
        \node[circle, draw=black!41, fill=black!22.55,
              minimum size=3.0mm, inner sep=0pt] (gc1) at (-0.40, -0.05) {};
        \node[circle, draw=black!41, fill=black!22.55,
              minimum size=3.0mm, inner sep=0pt] (gc2) at ( 0.40, -0.05) {};
        % Grandchildren
        \node[circle, draw=black!41, fill=black!16.4,
              minimum size=3.0mm, inner sep=0pt] (gg1) at (-0.55, -0.50) {};
        \node[circle, draw=black!41, fill=black!16.4,
              minimum size=3.0mm, inner sep=0pt] (gg2) at (-0.20, -0.50) {};
        \node[circle, draw=black!41, fill=black!16.4,
              minimum size=3.0mm, inner sep=0pt] (gg3) at ( 0.55, -0.50) {};
        % Edges (dashed, like a tree)
        \draw[treearrow] (gr) -- (gc1);
        \draw[treearrow] (gr) -- (gc2);
        \draw[treearrow] (gc1) -- (gg1);
        \draw[treearrow] (gc1) -- (gg2);
        \draw[treearrow] (gc2) -- (gg3);
    \end{scope}
    \node[anchor=north, align=center, font=\footnotesize] at (0,0)
        {Rule-based\\ derivation graphs\\ (M\O D pipeline)};
\end{scope}

% Arrow from MØD to Generate (with label)
\draw[grayA] (0.5, 12.10) -- (1.20, 12.10);

%% =================================================================
%%   STAGE D: LEGEND on right side
%% =================================================================
%\begin{scope}[shift={(11.5, 14.0)}]
%    % Bond/substructure proposals icon: small branched substructure
%    \begin{scope}[shift={(0,0)}]
%        % Bonds (drawn first, behind nodes)
%        \draw[SeaGreen, fig medium] (-0.35, 0.00) -- (-0.05, 0.00);
%        \draw[SeaGreen, fig medium] (-0.05, 0.00) -- ( 0.20, 0.22);
%        \draw[SeaGreen, fig medium] (-0.05, 0.00) -- ( 0.20,-0.22);
%        \draw[SeaGreen, fig medium] ( 0.20, 0.22) -- ( 0.20,-0.22);
%        % Atoms (small filled green dots with outline)
%        \foreach \x/\y in {-0.35/0.00, -0.05/0.00, 0.20/0.22, 0.20/-0.22} {
%            \draw[SeaGreen, fig line, fill=SeaGreen!55]
%                (\x,\y) circle (0.075);
%        }
%    \end{scope}
%    \node[anchor=west, font=\small, align=left] at (0.6, 0)
%        {bond /\\substructure\\proposals};
%
%    % Candidate fragments icon
%    \begin{scope}[shift={(0,-1.7)}]
%        \node[circle, fill=SeaGreen!50, draw=SeaGreen, minimum size=3mm, inner sep=0pt] (cf1) at (-0.2, 0.10) {};
%        \node[circle, fill=SeaGreen!25, draw=SeaGreen!70, minimum size=3mm, inner sep=0pt] (cf2) at ( 0.2, 0.05) {};
%        \node[circle, fill=SeaGreen!50, draw=SeaGreen, minimum size=3mm, inner sep=0pt] (cf3) at (-0.10,-0.30) {};
%        \node[circle, fill=white, draw=SeaGreen!70, minimum size=3mm, inner sep=0pt] (cf4) at ( 0.30,-0.30) {};
%        \draw[dashed, SeaGreen!70, thin] (cf1)--(cf2);
%        \draw[dashed, SeaGreen!70, thin] (cf3)--(cf4);
%    \end{scope}
%    \node[anchor=west, font=\small, align=left] at (0.6, -1.65)
%        {candidate\\fragments};
%
%    % Fragmentation DAG/tree icon
%    \begin{scope}[shift={(0,-3.4)}]
%        \node[fragroot, scale=0.8] (dr) at (0,0.30) {};
%        \node[fragnode, scale=0.8] (dl1) at (-0.30,-0.20) {};
%        \node[fragnode, scale=0.8] (dl2) at ( 0.30,-0.20) {};
%        \draw[treearrow] (dr)--(dl1);
%        \draw[treearrow] (dr)--(dl2);
%    \end{scope}
%    \node[anchor=west, font=\small, align=left] at (0.6, -3.30)
%        {fragmentation\\DAG / tree};
%\end{scope}
%
% =================================================================
%   STAGE E: Arrow Generate -> Score, plus return arrow (end-to-end)
% =================================================================
\draw[greenA] (5.0, 10.10) -- (5.0, 9.20);
\draw[blueA] (7.0, 9.20) -- (7.0, 10.10);
\node[font=\itshape\small, align=center] at (8.85, 9.65) {end-to-end\\[-0.7em] training};

% =================================================================
%   STAGE F: SCORE module (blue)
% =================================================================
\begin{scope}[shift={(6, 6.0)}]
    % Outer box
    \draw[rounded corners=3pt, RoyalBlue, fig medium, fill=RoyalBlue!10]
        (-4.8,-2.00) rectangle (4.8, 3.20);
    % Title
    \node[anchor=north, font=\LARGE\bfseries, text=RoyalBlue] at (0, 3.05) {Score};
    \node[anchor=north, font=\small\itshape] at (0, 2.55) {(neural embedding \& scoring module)};

    % Hexagonal graph icon (top-left)
    \begin{scope}[shift={(-3.4, 2.5)}]
        \foreach \i in {0,...,5} {
            \pgfmathsetmacro{\ang}{60*\i+30}
            \coordinate (h\i) at (\ang:0.32);
        }
        \draw[RoyalBlue, thick] (h0)--(h1)--(h2)--(h3)--(h4)--(h5)--cycle;
        \node[circle, fill=RoyalBlue, minimum size=1.4mm, inner sep=0pt] (m0) at (0,0) {};
        \foreach \i in {0,...,5} {
            \node[circle, fill=RoyalBlue, minimum size=1.4mm, inner sep=0pt] at (h\i) {};
            \draw[RoyalBlue, thin, ->, shorten >=0.6mm, shorten <=0.6mm]
                (m0) -- (h\i);
        }
    \end{scope}

    % Three sub-panels: Fragments, Molecule, Spectrum
    \draw[rounded corners=2pt, RoyalBlue!50, thin]
        (-4.5, 0.20) rectangle (4.5, 2.05);
    % Vertical separators
    \draw[RoyalBlue!40, thin] (-1.5, 0.20) -- (-1.5, 2.05);
    \draw[RoyalBlue!40, thin] ( 1.5, 0.20) -- ( 1.5, 2.05);

    \node[anchor=north, font=\small\bfseries] at (-3.0, 2.00) {Fragments};
    \node[anchor=north, font=\small\bfseries] at ( 0.0, 2.00) {Molecule};
    \node[anchor=north, font=\small\bfseries] at ( 3.0, 2.00) {Spectrum};

    % Fragments: a few small circles + dots
    \node[circle, fill=CornflowerBlue!60, draw=RoyalBlue, minimum size=3mm, inner sep=0pt] at (-3.7, 1.05) {};
    \node[circle, fill=CornflowerBlue!42, draw=RoyalBlue, minimum size=3mm, inner sep=0pt] at (-3.2, 1.10) {};
    \node[circle, fill=CornflowerBlue!54, draw=RoyalBlue, minimum size=3mm, inner sep=0pt] at (-2.7, 1.05) {};
    \node[circle, fill=CornflowerBlue!36, draw=RoyalBlue, minimum size=3mm, inner sep=0pt] at (-2.2, 1.10) {};
    \node[font=\small] at (-1.75, 1.05) {$\dots$};

    % Molecule: small graph
    \begin{scope}[shift={(0, 1.05)}]
        \node[circle, fill=NavyBlue, draw=NavyBlue, minimum size=3mm, inner sep=0pt] (mm1) at (-0.40, 0.10) {};
        \node[circle, fill=CornflowerBlue!60, draw=RoyalBlue, minimum size=3mm, inner sep=0pt] (mm2) at (-0.05, 0.20) {};
        \node[circle, fill=NavyBlue, draw=NavyBlue, minimum size=3mm, inner sep=0pt] (mm3) at ( 0.30, 0.05) {};
        \node[circle, fill=CornflowerBlue!60, draw=RoyalBlue, minimum size=3mm, inner sep=0pt] (mm4) at (-0.20,-0.25) {};
        \node[circle, fill=white, draw=RoyalBlue, minimum size=3mm, inner sep=0pt] (mm5) at ( 0.45,-0.30) {};
        \draw[thick] (mm1)--(mm2) (mm2)--(mm3) (mm2)--(mm4) (mm3)--(mm5);
        \draw[thick, dashed] (mm1)--(mm4);
    \end{scope}

    % Spectrum: vertical bars of varying heights
    \begin{scope}[shift={(2.0, 0.40)}]
        \draw[RoyalBlue!70, thin] (0,0) -- (2.2, 0);   % x axis
        \foreach \x/\h in {0.10/0.45, 0.30/0.20, 0.55/0.85, 0.75/0.55, 0.95/0.30,
                          1.25/0.95, 1.50/0.30, 1.70/0.65, 1.90/0.20, 2.10/0.10} {
            \draw[NavyBlue, fig medium] (\x, 0) -- (\x, \h);
        }
        \node[font=\small] at (2.4, 0.45) {$\dots$};
    \end{scope}

    % --- Shared latent space row ---
    \draw[rounded corners=2pt, RoyalBlue!50, thin, fill=RoyalBlue!6]
        (-4.5,-1.80) rectangle (4.5, 0.05);
    \node[anchor=west, font=\small\bfseries] at (-4.30,-0.875) {Shared latent space};

    % Scatter plot (right portion) — vertically centered in the box
    \begin{scope}[shift={(1.4,-0.90)}]
        \draw[->, black!70, thin] (-0.15,-0.80) -- (-0.15, 0.80);
        \draw[->, black!70, thin] (-0.15,-0.80) -- ( 2.65,-0.80);
        \foreach \px/\py in {
            -0.05/-0.30, 0.20/-0.10, 0.40/-0.45, 0.50/0.10, 0.65/-0.25,
            0.80/-0.05, 0.90/0.30, 1.05/0.00, 1.10/0.45, 1.20/-0.20,
            1.25/0.10, 1.35/-0.35, 1.40/0.20, 1.45/-0.05, 1.55/0.35,
            1.60/-0.15, 1.65/0.50, 1.75/0.05, 1.80/-0.30, 1.90/0.25,
            1.95/-0.05, 2.05/0.40, 2.10/-0.20, 2.20/0.15, 2.30/-0.30,
            0.30/0.55, 0.55/0.40, 1.00/-0.40, 1.50/0.55, 2.00/-0.50} {
            \node[circle, fill=RoyalBlue!70, draw=RoyalBlue, minimum size=1.4mm, inner sep=0pt] at (\px, \py) {};
        }
    \end{scope}

    % Arrow from Molecule to Shared latent space
    \draw[->, RoyalBlue, thick] (0, 0.50) -- (0, 0.10);
\end{scope}

% =================================================================
%   STAGE G: Two output boxes at the bottom
% =================================================================
% Spectrum prediction (left)
\begin{scope}[shift={(2.0, 2.55)}]
    \draw[rounded corners=3pt, black!60]
        (-2.4,-0.85) rectangle (2.4, 0.95);
    \node[anchor=south, font=\small\bfseries] at (0, 0.45) {Spectrum prediction};
    \begin{scope}[shift={(-2.0, -0.65)}]
        \draw[RoyalBlue!70, thin] (0,0) -- (4.1, 0);
        \foreach \x/\h in {0.15/0.55, 0.30/0.20, 0.50/0.40, 0.70/0.85, 0.90/0.30,
                          1.20/0.95, 1.45/0.45, 1.65/0.65, 1.90/0.30, 2.15/0.80,
                          2.40/0.50, 2.60/0.25, 2.85/0.70, 3.10/0.40, 3.35/0.85,
                          3.60/0.30, 3.85/0.55} {
            \draw[NavyBlue, fig medium] (\x, 0) -- (\x, \h);
        }
        \node[font=\small] at (4.30, 0.45) {$\dots$};
    \end{scope}
\end{scope}

% Inverse retrieval (right)
\begin{scope}[shift={(9.4, 2.55)}]
    \draw[rounded corners=3pt, black!60]
        (-2.7,-0.85) rectangle (2.7, 0.95);
    \node[anchor=south, font=\small\bfseries] at (0, 0.45) {Inverse retrieval};
    % Red spectrum
    \begin{scope}[shift={(-2.4, -0.65)}]
        \draw[BrickRed!70, thin] (0,0) -- (1.85, 0);
        \foreach \x/\h in {0.10/0.45, 0.25/0.85, 0.40/0.30, 0.55/0.65, 0.70/0.95,
                          0.90/0.40, 1.05/0.55, 1.20/0.20, 1.35/0.70, 1.50/0.30,
                          1.65/0.50, 1.78/0.25} {
            \draw[BrickRed, fig medium] (\x, 0) -- (\x, \h);
        }
    \end{scope}
    % Double-ended arrow
    \draw[<->, fig medium] (-0.30, -0.20) -- (0.30, -0.20);
    % Two small molecules + dots
    \begin{scope}[shift={(0.85,-0.30)}]
        \node[circle, fill=RoyalBlue, minimum size=2mm, inner sep=0pt] (im1) at (0,0.18) {};
        \node[circle, fill=white, draw=RoyalBlue, minimum size=2mm, inner sep=0pt] (im2) at (-0.20,-0.05) {};
        \node[circle, fill=CornflowerBlue!60, draw=RoyalBlue, minimum size=2mm, inner sep=0pt] (im3) at (0.20,-0.05) {};
        \node[circle, fill=NavyBlue, minimum size=2mm, inner sep=0pt] (im4) at (0.05,-0.30) {};
        \draw[thin] (im1)--(im2) (im1)--(im3) (im3)--(im4) (im2)--(im4);
    \end{scope}
    \begin{scope}[shift={(1.65,-0.30)}]
        \node[circle, fill=CornflowerBlue!60, draw=RoyalBlue, minimum size=2mm, inner sep=0pt] (jm1) at (0,0.18) {};
        \node[circle, fill=RoyalBlue, minimum size=2mm, inner sep=0pt] (jm2) at (-0.18,-0.10) {};
        \node[circle, fill=white, draw=RoyalBlue, minimum size=2mm, inner sep=0pt] (jm3) at (0.20,-0.05) {};
        \node[circle, fill=CornflowerBlue!60, draw=RoyalBlue, minimum size=2mm, inner sep=0pt] (jm4) at (0.05,-0.32) {};
        \draw[thin] (jm1)--(jm2) (jm1)--(jm3) (jm3)--(jm4);
    \end{scope}
    \node[font=\small] at (1.20, -0.50) {$\dots$};
    \node[font=\small] at (2.05, -0.50) {$\dots$};
\end{scope}

% Blue arrows from Score box to outputs
\draw[blueA] (3.0, 4.00) -- (2.5, 3.50);
\draw[blueA] (8.5, 4.00) -- (9.0, 3.50);

% =================================================================
%   STAGE H: (goal banner removed)
% =================================================================

\end{tikzpicture}%
    }
    \caption[Future extension toward a Generate -- Score architecture]{
        A future extension of the present work: shifting from fragment-structure \emph{encoding} to fragmentation-process \emph{learning},
        so that the fragmentation graph becomes a predicted object rather than a fixed input.
    }
    \label{fig:generative-future}
\end{figure}


\section{Advantages of Learning Fragment Generation}
Learning fragmentation generation would offer several advantages over the present setup.
First, it would reduce the dependence on exhaustive upstream preprocessing: instead of enumerating a full fragmentation space before learning begins, a learned generator could prioritize only the most informative or probable fragmentation events, thereby lowering the combinatorial burden and improving scalability.
Second, it would make the architecture more \emph{end-to-end}.
Where the spectrum predictor is currently conditioned on a fragment representation that is treated as given, an integrated design could optimize fragment proposals, fragment embeddings, and spectrum scoring jointly, so the model learns not only how to use fragments but which fragment structures are most helpful for downstream spectrum reconstruction and inverse retrieval.
Third, explicit fragmentation generation may improve interpretability.
Because the fragment structures are currently externally generated, the model cannot explain why one fragmentation path was preferred over another, whereas a learned generative stage could provide atom- or bond-wise probabilities over fragmentation events and thereby give a more direct account of the model's mechanistic assumptions.


\section{Scientific and Technical Challenges}
This transition is non-trivial and raises several challenges.
The first is supervision: spectral libraries generally provide molecule--spectrum pairs, but not ground-truth fragmentation graphs.
As in ICEBERG, the derivation graphs produced by the rule-based M\O D pipeline could serve as proxy labels for a learned generator, but such labels are not guaranteed to be chemically unique or
physically exact, and this ambiguity would need to be addressed carefully during training and evaluation.
The second is error propagation: once generation becomes learned, mistakes made by the generative stage may propagate into the spectral scoring stage.
This argues for staged training, beam search, uncertainty estimation, or latent-variable formulations that can maintain several plausible fragmentation hypotheses rather than committing to a single one.
The third is chemical realism: a learned generator should distinguish fragments that are structurally possible from those that are spectrometrically plausible, accounting for effects such as charge retention, neutral loss, and hydrogen rearrangements --- again favoring a modular Generate--Score design over a single black-box network.


\section{A Plausible Development Path}
A realistic path forward would proceed in stages.
The first stage would keep the current derivation graphs but train a neural module to rank or prune candidate fragmentation steps.
The second stage would learn to expand fragmentation graphs autoregressively, conditioned on the precursor graph and the partial fragmentation history.
The third stage would couple this generator tightly to the existing scoring model so that fragment generation, spectral prediction, and inverse retrieval are optimized jointly.
Seen in this way, the present system is not an alternative to Generate--Score approaches but an intermediate step toward them: it demonstrates that fragment-level intermediate representations can strengthen the connection between molecules and spectra, and the next major advance is to let the model learn which fragmentation structures to construct rather than only how to encode those supplied by an external pipeline.
This would bring the framework closer to a fully trainable, mechanistically informed, and end-to-end bidirectional model for structure--spectrum translation.
